\chapter{Calendario settimanale, Tab Ore e operazioni in massa}
\label{ch:calendario-ore}

\lettrine[lines=3]{N}{el corso} di aprile e maggio MMXXVI l'interfaccia
operativa di piTantum ha conosciuto la riscrittura pi\`u sostanziale
dalla nascita del progetto. Tre direttrici di lavoro si sono saldate
in un'unica esperienza editoriale: il \emph{calendario settimanale}
con drag-drop, la \emph{tab Ore} per la configurazione visuale degli
slot di lezione, e le \emph{operazioni in massa} sulla pagina Cattedre.
Il capitolo le racconta nell'ordine in cui le incontra l'utente che
parte dal modello dati e arriva alla griglia oraria.

\section{Il componente \texttt{WeeklyCalendarView}}
\label{sec:weekly-calendar}

\subsection{Origine e portata}

Fino al commit \texttt{c8075c0} la pagina \texttt{/schedule}
mostrava una matrice tabulare statica: righe = ore, colonne =
giorni, celle popolate da una stringa breve ``\textsc{materia}
$\cdot$ docente''. La modifica era possibile solo aprendo modali
o cambiando schede: nessun gesto diretto sulla cella, nessun
trascinamento. La nuova interfaccia ribalta il paradigma: un
\emph{unico} componente Svelte, \texttt{WeeklyCalendarView}, fa da
griglia condivisa per tutte le viste calendario del prodotto, in
modalit\`a \texttt{view}, \texttt{edit} e \texttt{schedule}.

\begin{defbox}[title={Il file \texttt{WeeklyCalendarView.svelte}}]
Una sola sorgente di verit\`a per la griglia settimanale.
Ricava i giorni e gli slot da \texttt{workingHoursStore}
(propagazione live, vedi \S\ref{sec:store-condiviso}), si
sintonizza al \emph{prop} \texttt{mode} e cambia comportamento
di conseguenza. Mille seicento righe di Svelte, ma una sola
gerarchia: \emph{slot} ${\rightarrow}$ \emph{evento} ${\rightarrow}$
\emph{azione}.
\end{defbox}

\subsection{Le tre modalit\`a}

\begin{description}
  \item[\textsc{view}] (sola lettura). Calendario di consultazione
    per docente, classe o aula: nessun drag, nessuna azione.
    \`E la modalit\`a usata, ad esempio, nelle pagine di
    diagnostica.
  \item[\textsc{edit}] (Tab Ore). La griglia non rappresenta
    lezioni ma \emph{slot disponibili}. L'utente disegna lo
    slot trascinando il mouse dall'alto in basso (cursore
    \texttt{grab}/\texttt{grabbing}), lo ridimensiona afferrando
    il bordo inferiore (cursore \texttt{ns-resize}), lo sposta
    afferrandone il corpo, lo cancella con un click.
  \item[\textsc{schedule}] (\texttt{/schedule}). La griglia
    contiene le lezioni effettive del problema risolto. Sono
    ammessi quattro gesti --- \emph{Modifica}, \emph{Sposta},
    \emph{Svincola}, \emph{Elimina} --- e il drag-drop
    consente sia il riposizionamento di una lezione gi\`a
    schedulata sia il piazzamento di una lezione \emph{non
    schedulata}, prelevata dal pool laterale.
\end{description}

\subsection{Rendering e interazioni}

\subsubsection{Anatomia dello slot}

Ogni cella della griglia \`e un \texttt{<div>} con classe
\texttt{cal-slot} e attributo \texttt{data-testid="sched-slot-\{D\}-\{H\}"},
dove \texttt{D} \`e l'indice di giorno (0$=$lun, $\dots$, 5$=$sab)
e \texttt{H} l'indice di ora. Sopra lo slot vivono le ``carte''
delle lezioni (classe \texttt{cal-event}, attributo
\texttt{data-testid="sched-lesson-\{id\}"}); ognuna porta tre micro
azioni in alto a destra: matita (modifica), freccia (svincola),
cestino (elimina). Tre cursori distinti accompagnano i tre stati
del drag (\texttt{grab} a riposo, \texttt{grabbing} mentre si
trascina, \texttt{ns-resize} sul bordo inferiore in modalit\`a
edit), come dettato dal commit \texttt{d4bc873}.

\subsubsection{La preview di conflitto}

Quando l'utente afferra una lezione e ne sorvola la griglia
(\texttt{dragover}), \texttt{WeeklyCalendarView} calcola in tempo
reale, su tutta la matrice settimanale, l'insieme delle celle
che genererebbero \emph{conflitto hard} con la lezione trascinata
(stesso docente in due posti, stessa classe in due posti, stessa
aula occupata). Le celle a rischio si tingono di rosso tenue
(classe CSS \texttt{cal-slot--drop-conflict}); l'evento sorgente
mostra un piccolo distintivo \emph{!}\ rosso
(\texttt{cal-event-conflict-badge}). \`E il \emph{soft-conflict
preview} introdotto dal commit \texttt{876755b}: l'utente
vede ci\`o che sta per fare \emph{prima} di farlo.

\subsection{Il pool delle lezioni non schedulate}
\label{sec:pool-unscheduled}

A destra della griglia, in modalit\`a schedule, una sidebar
\emph{pool} (id-test
\texttt{schedule-pool}, riga 1162 del file Svelte) elenca le
lezioni che esistono nel database ma non hanno
\texttt{day\_idx}/\texttt{hour\_idx}: tipicamente, lezioni
generate dall'ottimizzatore in stato sospeso o lezioni
\emph{svincolate} dall'utente con l'azione omonima (vedi
\S\ref{sec:azioni-quattro}). Ogni elemento del pool \`e
trascinabile e pu\`o essere droppato su qualsiasi slot per
schedularlo immediatamente.

\subsection{Le quattro azioni della modalit\`a schedule}
\label{sec:azioni-quattro}

\begin{description}
  \item[\textsc{Modifica}] (icona matita). Apre la modale di
    edit della singola \texttt{Lesson}: docente, materia, aula,
    classe, eventuale corteacher. Persiste tramite
    \texttt{PATCH /api/lessons/\{id\}}.
  \item[\textsc{Sposta}] (drag sulla cella di destinazione).
    Posta in \texttt{POST /api/lessons/\{id\}/move} con corpo
    \texttt{\{day\_idx, hour\_idx, on\_conflict\}}; il backend
    ritorna \texttt{200} se nessun conflitto, \texttt{409}
    altrimenti.
  \item[\textsc{Svincola}] (icona freccia). Toglie la lezione
    dalla griglia ma non la elimina dal database: la lezione
    finisce nel pool. Posta in \texttt{POST
    /api/lessons/\{id\}/unschedule}.
  \item[\textsc{Elimina}] (icona cestino, conferma richiesta).
    Cancella la \texttt{Lesson} dal database. Endpoint
    \texttt{DELETE /api/lessons/\{id\}}.
\end{description}

\section{La modale di conflitto: \emph{Sostituisci o Annulla}}
\label{sec:conflict-modal}

\subsection{Quando si attiva}

L'utente trascina una lezione (o un elemento del pool) su uno
slot \emph{gi\`a occupato} dal punto di vista del backend:
docente impegnato, classe impegnata, aula impegnata. Il
frontend prova prima un \emph{dry-run} (\texttt{POST
/api/lessons/\{id\}/move?on\_conflict=dry\_run}); il backend
risponde \texttt{409} con un payload del tipo
\begin{lstlisting}[basicstyle=\ttfamily\footnotesize]
{
  "detail": {
    "conflicts": {
      "teacher_busy":  [ {class_name, subject, teacher_name, ...}, ... ],
      "class_busy":    [ ... ],
      "room_busy":     [ ... ]
    }
  }
}
\end{lstlisting}
A questo punto si apre la modale \texttt{ScheduleConflictModal}
(commit \texttt{59ba4a8}).

\subsection{Le tre risposte e perch\'e nel drop sono solo due}

La modale \`e canonica: ovunque l'applicazione rilevi un
conflitto duro mostra lo stesso componente. Le risposte
possibili sono tre:
\begin{itemize}
  \item \textsc{Annulla}: chiude la modale senza azione, la
    lezione torna alla posizione di partenza.
  \item \textsc{Svincola}: rimuove dal calendario \emph{solo}
    le risorse confliggenti --- per i conflitti d'aula,
    azzera \texttt{classroom\_name}; per docente o classe,
    cancella la lezione confliggente. La lezione che si stava
    droppando si sistema dove voleva.
  \item \textsc{Elimina} / \textsc{Sostituisci}: cancella
    integralmente le righe \texttt{Lesson} confliggenti e
    piazza la nuova lezione al loro posto.
\end{itemize}
Nel flusso \emph{drop su slot occupato} di \texttt{/schedule}
l'opzione \textsc{Svincola} viene nascosta: piazzare una nuova
lezione \emph{senza} svincolare le altre porta inevitabilmente a
una sostituzione integrale, e la dicotomia \emph{Sostituisci /
Annulla} \`e quella che meglio descrive l'azione (prop
\texttt{showUnbind=false}, \texttt{deleteLabel="Sostituisci"}).

\begin{esempiobox}[title={Drop con conflitto, raccontato passo passo}]
\begin{enumerate}
  \item Trascino la lezione \texttt{2A\,$\cdot$\,Storia\,$\cdot$\,Rossi}
    da Luned\`{\i} 3 a Marted\`{\i} 4.
  \item Marted\`{\i} 4 \`e gi\`a occupato dalla lezione
    \texttt{2A\,$\cdot$\,Matematica\,$\cdot$\,Bianchi} in
    \texttt{Aula 12}.
  \item Drop: il frontend chiama
    \texttt{POST /api/lessons/17/move} con
    \texttt{on\_conflict=dry\_run}.
  \item Risposta 409: \texttt{class\_busy=[...]}.
  \item Si apre la modale: ``\textsc{conflitto rilevato},
    Marted\`{\i} 4 occupato da Bianchi.\hfill
    [\textsc{annulla}] [\textsc{sostituisci}]''.
  \item L'utente sceglie \textsc{sostituisci}.
  \item Il frontend ripete la chiamata con
    \texttt{on\_conflict=delete}: la lezione di Bianchi sparisce,
    la lezione di Rossi prende posto. La griglia ridisegna.
\end{enumerate}
\end{esempiobox}

\section{La Tab Ore (\texttt{/ore})}

\grokfig{shot_ore}{La Tab Ore: qui si definisce la settimana
scolastica (giorni e fasce orarie). Le modifiche si propagano in tempo
reale al calendario tramite \texttt{workingHoursStore}.}
\label{sec:tab-ore}

\subsection{Perch\'e una tab dedicata}

I codici scolastici italiani non hanno una griglia oraria
universale: alcune scuole sono su settimana corta (luned\`{\i}-
venerd\`{\i}), altre sul sabato; alcune fanno ore da \SI{60}{\minute},
altre da \SI{55} o \SI{50}{\minute}; alcune ammettono pause di
durata variabile. Hard-codare \texttt{DAYS = ['Lun', \dots, 'Sab']}
e \texttt{HOURS = 6}, come faceva l'engine fino a marzo MMXXVI,
significa pretendere che il mondo si conformi all'esempio.

Tre commit di refactoring hanno cambiato le carte in tavola:
\texttt{b24fdbd} (modello dati: tabelle \texttt{WorkingDay} e
\texttt{WorkingHourSlot}), \texttt{a3c12fb} (router REST), e
\texttt{bf6cfa2} (engine: lettura del config dalle nuove tabelle,
abbandono di \texttt{DAYS} e \texttt{HOURS} costanti). Sopra a
questo, il commit \texttt{b3236ab} ha introdotto la \emph{Tab
Ore} come pagina di prima classe in navbar.

\subsection{Modello dati e API}

\begin{itemize}
  \item \texttt{WorkingDay(id, dow, name, active)}: i giorni
    della settimana attivi nella scuola, con un \emph{name}
    leggibile (`luned\`{\i}', `mercoled\`{\i}', \dots).
  \item \texttt{WorkingHourSlot(id, day\_id, idx, label,
    start\_min, end\_min, kind)}: lo slot all'interno del giorno,
    con etichetta (1\textsuperscript{a} ora, 2\textsuperscript{a}
    ora, ricreazione, \dots), inizio e fine in minuti dalla
    mezzanotte, e \emph{kind} fra \textsc{lezione},
    \textsc{ricreazione}, \textsc{pausa}.
\end{itemize}

Il router REST \texttt{/api/working-hours/*} espone:
\begin{lstlisting}[basicstyle=\ttfamily\footnotesize]
GET    /api/working-hours/config             # bundle completo
GET    /api/working-hours/days
POST   /api/working-hours/days                # nuovo giorno
PUT    /api/working-hours/days/{day_id}
DELETE /api/working-hours/days/{day_id}       # cascade slot
PUT    /api/working-hours/days/{day_id}/slots # replace
POST   /api/working-hours/reset               # default lun-sab 8-14, 60'
\end{lstlisting}

L'endpoint \texttt{/reset} ripristina la configurazione di
default: \emph{luned\`{\i}-sabato, sei ore consecutive da
\SI{60}{\minute}, dalle 08:00 alle 14:00}.

\subsection{La UI vintage-style della Tab Ore}
\label{sec:ore-vintage}

I commit \texttt{d11e396} e \texttt{d4bc873} hanno trasformato
l'editor degli slot in una scheda di lavoro tipograficamente
gradevole, ispirata --- con qualche libert\`a --- ai registri
manoscritti di scuola del primo Novecento: caratteri serif sobri,
nessun pulsante baroccamente colorato, un solo accento di colore
(indaco) per le azioni, e tre cursori CSS distinti per
trasmettere il comportamento del drag senza spiegarlo a parole.

\begin{itemize}
  \item Su ogni cella di slot esistente compaiono due bottoncini
    sobri: \emph{Modifica} (apre un popover inline che si apre
    \emph{sopra} la cella, mai sotto, per non sporcare il flusso
    di lettura) e \emph{Cancella} (con doppio click di
    conferma).
  \item Le \emph{drop key hints} guidano l'utente: trascinando
    sul bordo inferiore della cella si \emph{ridimensiona} lo
    slot; trascinando dal corpo lo si \emph{sposta}. Le hint
    appaiono solo durante il drag.
  \item Il \emph{fantasmino traslucido}: quando si crea un nuovo
    slot trascinando il mouse, una preview semitrasparente segue
    il cursore. Questo \`e ci\`o che il commit \texttt{d4bc873}
    chiama \emph{ghost element}.
\end{itemize}

\subsection{Il fix che ha sbloccato la suite Cypress}

Una nota a margine, ma istruttiva. Il commit \texttt{5c1ba34}
porta il messaggio ``\emph{namespace-import was masking api.get}''.
Tradotto: il file \texttt{ore/+page.svelte} importava
\texttt{api} come namespace (\texttt{import * as api}), ma una
variabile locale \texttt{api} dichiarata pi\`u in basso oscurava
il namespace, rompendo \texttt{api.get(...)} a runtime. La
correzione (\texttt{import \{ api \} from \dots}) ha portato la
suite Cypress della Tab Ore da \texttt{0/15} a \texttt{15/15}
verde. La lezione: il namespace import \`e fragile in presenza
di variabili omonime; meglio l'import nominale.

\section{Lo store condiviso \texttt{workingHoursStore}}
\label{sec:store-condiviso}

\subsection{Il problema}

Cambiando uno slot in \texttt{/ore}, fino al commit
\texttt{e0b76bc} le altre viste calendario non si accorgevano
del cambiamento finch\'e non si ricaricava la pagina. Esempio
tipico: aggiungo la 7\textsuperscript{a} ora nella Tab Ore, salgo
in \texttt{/schedule} e la griglia mostra ancora sei colonne.

\subsection{La soluzione}

\texttt{workingHoursStore} \`e uno store Svelte (Writable) che
conserva il bundle di \texttt{WorkingDay} e \texttt{WorkingHourSlot}
e si auto-aggiorna ogni volta che la Tab Ore conferma un save.
Tutti i \texttt{WeeklyCalendarView} sottoscritti allo store
ricevono il nuovo array di slot e ridisegnano. Niente reload,
niente prop-drilling: la propagazione \`e \emph{live}.

\section{Operazioni in massa su \texttt{/assignments}}
\label{sec:bulk-assignments}

I commit \texttt{e6314eb} (backend) e \texttt{16f882a} (frontend)
hanno completato la pagina \emph{Cattedre} con il pannello
\emph{operazioni in massa}. La selezione multipla \`e gi\`a presente
in altre pagine ma qui apre il ventaglio di cinque azioni:

\begin{itemize}
  \item \textsc{Elimina}: cancella le righe selezionate
    (\texttt{POST /api/assignments/bulk-delete}).
  \item \textsc{Lock}/\textsc{Unlock}: alza o abbassa il flag
    \texttt{locked} (\texttt{POST /api/assignments/bulk-set-flag}
    con \texttt{flag=locked}).
  \item \textsc{Cambia docente}: assegna un nuovo docente in
    blocco (\texttt{POST /api/assignments/bulk-change-teacher}).
  \item \textsc{Imposta flag}: una versione generica di
    \emph{lock} per i flag arbitrari supportati dallo schema.
\end{itemize}

L'utente seleziona le righe (Ctrl-click, Shift-click, oppure il
checkbox di intestazione per ``seleziona tutto''), apre il
pannello e conferma; il frontend manda \emph{una sola} richiesta
al backend, che apre una transazione e applica l'operazione in
blocco. Per cattedre numerose --- una scuola di duecento
cattedre non \`e rara --- la differenza rispetto a duecento
chiamate seriali \`e di un ordine di grandezza.

\section{Il quadro d'insieme}

Le quattro innovazioni descritte in questo capitolo non sono
indipendenti: si saldano in un'unica esperienza editoriale.

\fleuron

La Tab Ore configura la griglia. Lo store condiviso la propaga.
\texttt{WeeklyCalendarView} la rende. La modale di conflitto
risolve l'ultima ambiguit\`a quando l'utente compone, smonta e
ricompone la sua settimana scolastica con un dito.
Nelle \emph{piccole cose} dell'interazione --- un cursore
\texttt{ns-resize} sul bordo, una preview rossa quando il drop
sarebbe scorretto, un fleuron al posto di un bullet quadrato ---
si misura la differenza fra una pagina web e un manuale di
scuola.
